\subsection{支配特征标}

将 \(X_{+}\) 定义为如下集合：其元素为 \(\lambda\)，且属于 \(\mathrm{X(T)}\)，并满足 \(\langle\lambda,x\rangle\geq0\) 对所有 \(x\in\Gamma(\mathrm{T})_{++}\) 成立；也就是线性形式 \(\delta(\lambda):\mathfrak{t}_{\mathbf{C}}\to\mathbf{C}\) 将 t 的房 C 映到 \(iR_{+}\)。

赋予 X(T) 以序群结构，其中正元素是 \(X_{+}\) 中的元素；置 \(R_{+} = R(G, T) \cap X_{+}\)，\(R_{-} = -R_{+}\)。\(R_{+}\) 的元素称为正根，\(R_{-}\) 的元素称为负根；每个根或者是正根，或者是负根（第 VI 章，第 1 节，第 6 号，定理 3）。不可能写成两个正根之和的正根称为单根；每个正根都是单根之和（同上）；单根构成 X(T) 中由根生成的子群的一个基，该子群可与 X(T/C(G)) 识别（第 4 节，第 4 号）；关于单根的反射生成 Weyl 群 \(W = W_{G}(T)\)（第 VI 章，第 1 节，第 5 号，定理 2）。

引理 1. 令 \(\lambda\) 为 \(\mathrm{X(T)}\) 的一个元素。以下条件等价：

\begin{enumerate}
\def\labelenumi{(\roman{enumi})}
\item
  \(\lambda - w(\lambda) \geq 0\)（分别为 \(>0\)）对所有 \(w \in \mathrm{W}\) 满足 \(w \neq 1\) 成立；
\item
  当 \(w \in \mathrm{W}\) 且 \(w \neq 1\) 时，\(\lambda - w(\lambda)\) 是单根的具有正系数的线性组合（分别为不全为零的正系数线性组合）；
\item
  \(\langle \lambda, K_{\alpha} \rangle \geq 0\)（分别为 \(>0\)）对每个正根 \(\alpha\) 成立；
\item
  \(\langle\lambda,K_{\alpha}\rangle\geq0\)（分别为 >0）对每个单根 \(\alpha\) 成立。
\end{enumerate}

条件（iii）与（iv）的等价性是立即的。由于 \(K_{\alpha}\) 的集合可与 \(\mathrm{R}(\mathrm{G},\mathrm{T})\) 的逆根系识别（第 4 节，第 5 号），（i）与（iii）的等价性由第 VI 章第 1 节第 6 号命题 18 及其推论得到。(ii) \(\Rightarrow\) (i) 是平凡的，反向蕴含由同上引文得到。

将 \(X_{++}\) 记为 \(\mathrm{X(T)}\) 中满足 \(\langle\lambda,K_{\alpha}\rangle\geq0\) 对每个正根 \(\alpha\) 成立的元素集合。\(X_{++}\) 的元素称为支配的。它们构成 W 在 \(\mathrm{X(T)}\) 上作用的基本域（第 VI 章，第 1 节，第 10 号）。有 \(X_{++}\subset X_{+}\)。

若 G 单连通，则对于每个单根 \(\alpha\)，存在 X(T) 的元素 \(\varpi_{\alpha}\)，使得 \(\langle\varpi_{\beta}, K_{\alpha}\rangle = \delta_{\alpha\beta}\) 对每个单根 \(\beta\) 成立，即 \(s_{\alpha}(\varpi_{\alpha}) = \varpi_{\alpha} - \alpha\)，且 \(s_{\beta}(\varpi_{\alpha}) = \varpi_{\alpha}\) 对每个单根 \(\beta \neq \alpha\) 成立；\(\varpi_{\alpha}\) 称为基本支配权；它们构成交换群 \(\mathrm{X(T)}\) 和交换幺半群 \(\mathrm{X}_{++}\) 的基；更精确地，每个元素 \(\lambda\) 属于 \(\mathrm{X(T)}\)，都可写成 \(\lambda = \sum_{\alpha} \langle \lambda, K_{\alpha} \rangle \varpi_{\alpha}\) 的形式。

将 \(\rho\) 记为 \(\mathrm{X(T)}\otimes \mathbf{Q}\) 中满足下式的元素：

\[
2 \rho = \sum_ {\alpha \in \mathrm{R} _ {+}} \alpha .
\]

\(\langle\rho,K_{\alpha}\rangle=1\) 对每个单根 \(\alpha\) 成立（第 VI 章，第 1 节，第 10 号，命题 29）。若 G 单连通，则 \(\rho\) 是基本支配权之和。

